\documentclass[12pt]{article}

\usepackage{sbc-template}
\usepackage{graphicx,url}
\usepackage[utf8]{inputenc}
\usepackage[T1]{fontenc}
\usepackage[brazil]{babel}
\usepackage{booktabs}
\usepackage{amsmath}

\graphicspath{{figuras/}}

\sloppy

\title{O que um Mixture of Experts ganha sobre o modelo denso?\\
Um estudo com GPT-2 treinado do zero em português}

\author{Raisson Adrian Grangeiro Souto\inst{1}}

\address{Programa de Pós-Graduação em Ciência da Computação (PPGCC)\\
  Universidade Federal de Campina Grande (UFCG) -- Campina Grande, PB -- Brasil
  \email{raisson.souto@copin.ufcg.edu.br}
}

\begin{document}

\maketitle

\begin{abstract}
  % TODO: abstract em inglês (máx. 10 linhas). Escrever por último.
\end{abstract}

\begin{resumo}
  % TODO: resumo em português (máx. 10 linhas). Escrever por último.
\end{resumo}

% Limite: 8 páginas no modelo SBC.
% Rubrica (Nota B): motivação/hipóteses/relacionados 2,0; desenho
% experimental 2,5; resultados e análise 3,0; limitações 1,0; escrita 1,5.

\section{Introdução}\label{sec:intro}
% Contexto: MoE como forma de aumentar parâmetros sem aumentar o custo por
%   token \cite{shazeer2017moe,fedus2022switch}.
% Problema: em escala reduzida (~10M parâmetros, ~27M tokens), o ganho do
%   MoE sobre o denso não é óbvio.
% Pergunta central: o que um MoE ganha sobre o modelo denso, na mesma
%   escala e com o mesmo orçamento?
% Contribuições e organização do artigo.

\section{Hipóteses}\label{sec:hipoteses}
% Escritas ANTES dos experimentos (exigência da rubrica).
% H1: MoE com mesmos parâmetros ativos (treino 2) atinge loss de validação
%     menor que o vanilla com o mesmo orçamento de tokens.
% H2: MoE com mesmos parâmetros totais (treino 3) fica próximo do vanilla
%     com menor custo de FFN por token.
% H3: sem a loss de balanceamento, a utilização dos especialistas colapsa.
% H4: especialistas se especializam por tipo de token (funcionais, números,
%     citações, pontuação) mais do que por área (sigla_centro).
% H5: hipótese da modificação escolhida (treino 4).

\section{Trabalhos Relacionados}\label{sec:relacionados}
% Citar os artigos originais de cada técnica implementada.

\subsection{Transformers decoder-only}
% \cite{vaswani2017attention}, \cite{radford2019gpt2}, \cite{raschka2024build}

\subsection{Mixture of Experts}
% \cite{shazeer2017moe}, \cite{lepikhin2021gshard}, \cite{fedus2022switch},
% \cite{jiang2024mixtral}, \cite{qwen3}

\subsection{Eficiência de treino}
% \cite{micikevicius2018mixed}, \cite{dao2022flashattention},
% \cite{loshchilov2019adamw}

\subsection{Modificação escolhida}
% Conforme o treino 4: \cite{su2024roformer} (RoPE),
% \cite{ainslie2023gqa} (GQA), \cite{zhang2019rmsnorm} e
% \cite{shazeer2020glu} (RMSNorm + SwiGLU).

\section{Metodologia}\label{sec:metodologia}

\subsection{Corpus e tokenização}
% Introduções de 23.026 trabalhos da UFCG, 11 centros; partição estratificada
% (treino 20.722 docs / ~27M tokens; validação 1.152 / ~1,6M; teste após o
% prazo). Tokenizador BPE de 16k de referência \cite{sennrich2016bpe}.

\subsection{Modelo base e treino otimizado}
% GPT-2 do livro: 4 camadas, 4 cabeças, emb 256, contexto 256, vocab 16k.
% Loop do Apêndice D: warmup linear, cosine decay, clipping, AdamW.
% Otimizações: fp16 + GradScaler, SDPA/FlashAttention, AdamW fused,
% torch.compile. Tab.~\ref{tab:otimizacoes}.

\subsection{Camada MoE}
% Roteador linear, top-k (8 especialistas, top-2), loss auxiliar de
% balanceamento estilo Switch \cite{fedus2022switch}:
% \begin{equation}
%   \mathcal{L}_{aux} = \alpha \, N \sum_{i=1}^{N} f_i \, P_i
% \end{equation}
% Dimensionamento dos treinos 2 (mesmos parâmetros ativos) e 3 (mesmos
% parâmetros totais). Tab.~\ref{tab:configs}.

\subsection{Protocolo experimental}
% Mesmo orçamento de tokens, mesmo contexto, mesmo cronograma de LR,
% semente fixa (123), uma rodada por configuração.
% Métricas: loss/perplexidade e BPB na validação, parâmetros totais e
% ativos, tokens/s, memória de pico, tempo de parede, amostras para os 3
% prompts fixos, utilização dos especialistas. BPB no teste após o prazo.
% Logs versionados; hash SHA-256 dos checkpoints no README.

\begin{table}[ht]
\centering
\caption{Configurações comparadas.}\label{tab:configs}
\begin{tabular}{@{}lrrrr@{}}
\toprule
Treino & Especialistas & Top-$k$ & Parâm.\ totais & Parâm.\ ativos \\
\midrule
1 -- Vanilla        & -- & -- & & \\
2 -- MoE (ativos)   & 8  & 2  & & \\
3 -- MoE (totais)   & 8  & 2  & & \\
4 -- Modificação    & -- & -- & & \\
\bottomrule
\end{tabular}
\end{table}

\section{Resultados}\label{sec:resultados}

\subsection{Treino otimizado}
% Tokens/s e memória de pico antes e depois de cada otimização (~200 passos).

\begin{table}[ht]
\centering
\caption{Efeito de cada otimização no throughput e na memória.}\label{tab:otimizacoes}
\begin{tabular}{@{}lrr@{}}
\toprule
Configuração & Tokens/s & Memória de pico (GB) \\
\midrule
Base (fp32)            & & \\
+ fp16 e GradScaler    & & \\
+ FlashAttention (SDPA)& & \\
+ AdamW fused          & & \\
+ torch.compile        & & \\
\bottomrule
\end{tabular}
\end{table}

\subsection{Vanilla vs.\ MoE}
% Curvas de loss de treino/validação (Fig.); perplexidade e BPB finais.
% Comparação por parâmetros ativos vs. por parâmetros totais.

\subsection{Utilização dos especialistas}
% Utilização ao longo do treino (há colapso?).
% Que tokens cada especialista recebe: funcionais, números, citações
% (AUTOR, ano), pontuação.
% Especialização por área (sigla_centro).

\subsection{Modificação escolhida}
% Treino 4 vs. vanilla.

\subsection{Amostras geradas}
% Prompts: ``Nos últimos anos,'', ``Este trabalho tem como objetivo'',
% ``No contexto brasileiro,''.

\section{Discussão}\label{sec:discussao}
% Retomar H1--H5. Explicar por que funcionou ou não. Cautela com
% diferenças pequenas (uma semente por configuração).

\section{Limitações e Ameaças à Validade}\label{sec:limitacoes}
% Escala reduzida, uma semente, uma época, GPU T4 (sem bf16), corpus
% de domínio único.

\section{Conclusão}\label{sec:conclusao}
% Resposta à pergunta central; trabalhos futuros.

\bibliographystyle{sbc}
\bibliography{references}

\end{document}
